\section{Препроцессинг данных}

Проводится два эксперимента с разными условиями подготовки данных. Все изменения касаются только стороны данных — архитектуры и гиперпараметры одинаковы в обоих.

\subsection{Общий пайплайн (оба эксперимента)}

\subsubsection{Схема разметки BIO}

Сущности кодируются в формате BIO: \texttt{B-X} --- первый токен сущности типа X, \texttt{I-X} --- продолжение, \texttt{O} --- вне сущности. Итого 13 меток.

Альтернатива --- схема BIOES (добавляет \texttt{S-X} для одиночных и \texttt{E-X} для последних токенов). BIOES точнее кодирует границы и теоретически упрощает задачу для CRF, однако увеличивает пространство меток с 13 до 21. На практике для коротких span'ов, типичных для SciERC, выигрыша не наблюдается, поэтому выбран стандартный BIO.

\subsubsection{Токенизация и выравнивание меток}

WordPiece (BERT, SciBERT) и BPE (RoBERTa) разбивают слова на субтокены. Возникает вопрос: какую метку присваивать продолжающим субтокенам?

\textbf{Вариант «-100»:} продолжающие субтокены маскируются --- модель учится только на первом субтокене каждого слова. Это стандартный подход в ранних реализациях. Недостаток: большая часть обучающего сигнала теряется; для длинных слов, характерных для научных текстов (\textit{convolutional}, \textit{bidirectional}), модель видит только первый фрагмент.

\textbf{Вариант «I-X» (используется в обоих экспериментах):} продолжающие субтокены получают метку \texttt{I-X}, если слово начинается с \texttt{B-X}, или ту же \texttt{I-X}, если слово уже было \texttt{I-X}. Это увеличивает количество обучающих примеров и делает предсказания более согласованными на уровне субтокенов. Трейдоф: для однотокенных сущностей, разбитых на несколько субтокенов, модель вынуждена предсказывать \texttt{B-X} на первом и \texttt{I-X} на остальных --- технически правильно, но создаёт слабо мотивированный \texttt{I-X} без предшествующего полного слова.

\subsubsection{Динамический паддинг}

Паддинг выполняется не при токенизации, а при формировании батча: последовательности дополняются до длины самой длинной в батче. Преимущество: при типичных длинах предложений SciERC (20--60 токенов) паддинг до 512 означал бы, что 85--90\% вычислений атеншена тратится на паддинговые токены. Динамический паддинг устраняет эту неэффективность без потери качества.

\subsection{Эксперимент 1 --- базовый препроцессинг}

Каждое предложение обрабатывается независимо, аугментация и межпредложенческий контекст не применяются.

\subsection{Эксперимент 2 --- улучшенный препроцессинг}

\subsubsection{Аугментация заменой сущностей (Entity Swap)}

\textbf{Мотивация.} BERT-based NER-модели склонны запоминать поверхностные формы сущностей: если \textit{LSTM} встречается в обучающей выборке как Method, модель может учиться на имени, а не на контексте. На маленьком датасете (350 документов) это особенно опасно. Entity swap нарушает эту корреляцию: заменяя сущность на другую того же типа, мы сохраняем синтаксический контекст, но меняем поверхностную форму, вынуждая модель опираться на окружение.

\textbf{Реализация.} По обучающей выборке строится словарь: для каждого типа сущностей собираются все встречающиеся span'ы. При аугментации каждый span в предложении заменяется с вероятностью $p = 0.5$ случайным span'ом того же типа. Замена только внутри типа принципиальна: межтиповая замена порождала бы противоречивые примеры (Method на месте Metric в том же синтаксическом контексте).

\textbf{Трейдофы.}
\begin{itemize}
    \item Аугментированные предложения семантически могут быть неестественными (редкая метрика в контексте, типичном для метода). Это шум, но регуляризирующий.
    \item Датасет увеличивается до $\sim 2 \times$ --- обучение занимает вдвое больше времени при том же числе эпох.
    \item Для очень редких типов (Metric, Task) словарь мал, и разнообразие замен ограничено.
\end{itemize}

\subsubsection{Межпредложенческий контекст}

\textbf{Мотивация.} BERT кодирует каждое предложение независимо, не видя соседних. Однако в научных текстах именованные сущности часто вводятся и используются повторно в нескольких предложениях подряд. Граничные токены первого и последнего слова предложения лишены левого и правого контекста соответственно.

\textbf{Реализация.} Для каждого предложения $s_i$ в документе к нему приписываются соседи с окном $w = 1$:
\[
    \underbrace{s_{i-1}}_{\text{контекст, метка} = -100} \;\; \underbrace{s_i}_{\text{цель}} \;\; \underbrace{s_{i+1}}_{\text{контекст, метка} = -100}
\]
Токены контекстных предложений получают маску $-100$: функция потерь не вычисляется по ним, но энкодер обрабатывает их наравне с целевыми. Контекст применяется к train, dev и test: это корректно, поскольку на инференсе соседние предложения также доступны.

\textbf{Трейдофы.}
\begin{itemize}
    \item Длина входной последовательности увеличивается примерно втрое. При лимите 512 токенов предложения длиннее $\sim 170$ токенов с контекстом будут усечены. В SciERC такие случаи единичны.
    \item Окно $w > 1$ даёт больше контекста, но при $w = 2$ риск усечения центрального предложения существенно возрастает, а выигрыш от дальних предложений снижается.
    \item Вычислительные затраты растут квадратично по длине последовательности из-за механизма внимания.
\end{itemize}

\subsection{Взвешивание классов}

\textbf{Мотивация.} В SciERC метка \texttt{O} составляет около 70\% всех токенов, тогда как Metric --- около 2\%. При стандартной кросс-энтропии модель может игнорировать редкие классы и получать низкий лосс, просто предсказывая \texttt{O} везде.

\textbf{Обратная частота} $w_c = N / (K \cdot n_c)$ исправляет это, но создаёт соотношение весов до $35 \times$ между \texttt{O} и Metric --- модель начинает чрезмерно «видеть» сущности, снижая точность.

\textbf{Sqrt-сглаживание:}
\[
    w_c = \sqrt{\frac{N}{K \cdot n_c}}
\]
уменьшает соотношение примерно до $6 \times$, балансируя между полным игнорированием дисбаланса и его гиперкоррекцией. Квадратный корень --- эвристика без строгого обоснования; альтернатива --- настройка коэффициентов вручную, что нецелесообразно при 13 классах.

Каждая модель обучается в вариантах с весами и без, чтобы эмпирически проверить, помогает ли взвешивание именно на этом датасете.

\subsection{Метрика оценки}

Качество оценивается по entity-level macro F1 (библиотека \texttt{seqeval}). Сущность считается правильно распознанной только при точном совпадении границ span'а \textit{и} типа. Macro F1 усредняет F1 по всем 6 типам с равными весами, что делает метрику чувствительной к редким классам --- именно здесь сосредоточена практическая сложность задачи.

Альтернатива --- micro F1 --- взвешивает по числу сущностей каждого типа, что даёт непропорционально большой вес частым типам (Generic, Method) и скрывает плохое качество на Metric и Task.
